% =============================================================================
\section{Un topic propio}
% =============================================================================

\begin{frame}{battery\_publisher y battery\_subscriber}
  \centering
  \begin{tikzpicture}[every node/.style={font=\scriptsize},
      bx/.style={draw=azul,fill=azulclaro,rounded corners,minimum height=0.8cm,minimum width=2.4cm,align=center}]
    \node[bx,fill=gray!15,draw=gray] (px) at (0,0) {PX4};
    \node[bx] (pub) at (4.4,0) {battery\_publisher\\(traduce)};
    \node[bx] (sub) at (9.4,0) {battery\_subscriber\\(muestra)};
    \node[bx,fill=gray!10,draw=gray] (any) at (9.4,-1.6) {cualquier otro nodo\\o \cmd{ros2 topic echo}};
    \draw[-{Stealth}] (px) -- node[above,font=\tiny,align=center]{battery\_status\_v1\\BatteryStatus} (pub);
    \draw[-{Stealth},thick,naranja] (pub) -- node[above,font=\tiny,align=center]{/vespo/battery\_percent\\std\_msgs/Float32} (sub);
    \draw[-{Stealth},thick,naranja] (pub.south) |- (any.west);
  \end{tikzpicture}
  \vspace{0.6em}
  \begin{itemize}\small
    \item Un nodo \textbf{traduce} lo que dice PX4 a algo simple: solo el porcentaje.
    \item Cualquier otro nodo lo usa \textbf{sin saber nada} de \cmd{px4\_msgs}: una interfaz, la misión, un
      registro\ldots
    \item Es el patrón para pasarse datos entre nodos nuestros.
  \end{itemize}
\end{frame}

\begin{frame}[fragile]{El publisher}
  \begin{lstlisting}[style=cpp,caption={battery\_publisher.cpp}]
// Publisher con QoS por defecto: rclcpp::QoS(10) = reliable, cola de 10.
percent_pub_ = create_publisher<Float32>("/vespo/battery_percent", 10);

auto px4_qos = rclcpp::QoS(1).best_effort().transient_local();
battery_sub_ = create_subscription<BatteryStatus>(
  "/fmu/out/battery_status_v1", px4_qos,
  [this](const BatteryStatus::SharedPtr msg) {
    // remaining = -1 cuando PX4 aún no sabe: no publicamos basura.
    if (!msg->connected || msg->remaining < 0.0f) {
      return;
    }
    Float32 out;
    out.data = msg->remaining * 100.0f;   // 0..1 -> 0..100 %
    percent_pub_->publish(out);
  });
  \end{lstlisting}
  \small Publica \textbf{dentro} del callback: cada vez que llega una batería de PX4 (1 Hz), sale un
  porcentaje. No hace falta timer.
\end{frame}

\begin{frame}{Las decisiones de diseño}
  \centering\small
  \begin{tabular}{lp{7.6cm}}
    \toprule
    Decisión & Por qué \\
    \midrule
    Nombre \cmd{/vespo/...} & Convención del equipo: nuestros topics no se confunden con los de PX4
      (\cmd{/fmu/...}) ni de otros paquetes. \\
    \cmd{std\_msgs/Float32} & Un mensaje estándar con un solo campo (\cmd{data}): no hace falta crear uno. \\
    QoS por defecto (\cmd{10}) & Entre nodos de ROS 2 no hace falta el best\_effort de PX4: reliable, cola de 10. \\
    Filtrar \cmd{remaining < 0} & $-1$ = ``PX4 no sabe'': mejor no publicar que publicar basura. \\
    \bottomrule
  \end{tabular}
\end{frame}

\begin{frame}[fragile]{El subscriber}
  \begin{lstlisting}[style=cpp,caption={battery\_subscriber.cpp}]
percent_sub_ = create_subscription<Float32>(
  "/vespo/battery_percent", 10,
  [this](const Float32::SharedPtr msg) {on_percent(msg->data);});

void on_percent(float percent) {
  // Barra de 20 caracteres: cada # vale 5 %.
  const int filled = static_cast<int>(percent / 5.0f + 0.5f);
  std::string bar;
  for (int i = 0; i < 20; i++) {
    bar += i < filled ? "#" : ".";
  }
  if (percent < low_percent_) {
    RCLCPP_WARN(get_logger(), "[%s] %5.1f %%  BAJA", bar.c_str(), percent);
  } else {
    RCLCPP_INFO(get_logger(), "[%s] %5.1f %%", bar.c_str(), percent);
  }
}
  \end{lstlisting}
\end{frame}

\begin{frame}{Detalles del subscriber}
  \begin{columns}[c]
    \column{0.52\textwidth}\small
    \begin{itemize}
      \item El callback llama a un método (\cmd{on\_percent}) con \textbf{solo el número}: el resto del código
        no sabe de mensajes.
      \item \cmd{+ 0.5f} antes de convertir a \cmd{int} \textbf{redondea} (sin él, \cmd{static\_cast} trunca:
        9.9 $\rightarrow$ 9).
      \item \cmd{bar.c\_str()}: \cmd{\%s} pide texto de C (\cmd{const char*}); \cmd{c\_str()} lo saca de un
        \cmd{std::string}.
    \end{itemize}
    \column{0.45\textwidth}
    \begin{terminal}{ros2 run px4\_examples battery\_subscriber}
[\#\#\#\#\#\#\#\#\#\#\#\#\#\#\#\#\#\#\#\#] 100.0 \%\\
[\#\#\#\#\#\#\#\#\#\#\#\#\#\#\#\#\#.\,.\,.] \ 84.5 \%\\
[\#\#\#\#\#.\,.\,.\,.\,.\,.\,.\,.\,.\,.\,.\,.\,.\,.\,.] \ 25.0 \%\ \ BAJA
    \end{terminal}
  \end{columns}
\end{frame}

\begin{frame}[fragile]{Probarlo sin PX4}
  \begin{lstlisting}[style=bash]
ros2 run px4_examples battery_subscriber            # terminal 1
ros2 topic pub /vespo/battery_percent std_msgs/msg/Float32 "{data: 15.0}"   # terminal 2
ros2 topic echo /vespo/battery_percent              # ver lo que pasa por el topic
ros2 topic info /vespo/battery_percent              # quién publica y quién escucha
  \end{lstlisting}
  \begin{idea}
    \small Como el subscriber solo depende del topic, se puede probar \textbf{sin simulación}: publiquen
    valores a mano para ver el aviso de batería baja. Así se prueban las partes de una misión por separado.
  \end{idea}
\end{frame}

\begin{frame}{Ejercicios}
  \begin{enumerate}\small
    \item Agreguen al publisher un segundo topic \cmd{/vespo/battery\_voltage} con el voltaje.
    \item Hagan que el subscriber avise \textbf{una sola vez} al cruzar el umbral (no en cada mensaje).
    \item Publiquen la altura actual (\cmd{-z} de \cmd{vehicle\_local\_position\_v1}) en
      \cmd{/vespo/altitude} y grafíquenla con \cmd{rqt\_plot}.
    \item Cambien el QoS del subscriber a \cmd{best\_effort}: ¿sigue funcionando? ¿Y si el que cambia es el
      publisher?
  \end{enumerate}
\end{frame}
